\section*{Sommario}
\addcontentsline{toc}{section}{Sommario}

Questa relazione documenta la riproduzione, in OpenFOAM-13 e con la libreria
termochimica \mpp, dei casi di verifica zero-dimensionali con cui Casseau
\textit{et al.}~\cite{casseau2016} hanno validato il solver a due temperature
\hyfoam. Il modello di Park a due temperature -- una temperatura
traslazionale-rotazionale $\Ttr$ e una vibro-elettronica $\Tve$ -- è stato
implementato in forma conservativa, risolvendo le equazioni nelle energie
anziché nelle temperature e delegando a \mpp{} sia l'inversione delle energie
in temperature sia il calcolo delle proprietà termodinamiche.

Il lavoro procede su due percorsi che condividono le stesse funzioni sorgente:
un insieme di programmi zero-dimensionali autonomi, che integrano le equazioni
del bagno termico chiamando soltanto \mpp{} e coprono tutte le figure del
capitolo di verifica del paper, e il solver CFD, che risolve le stesse equazioni su una singola cella
attraverso un modello termodinamico che fa da ponte verso la libreria. L'accordo fra i due percorsi, entro lo 0{,}03--0{,}2\%
su tutte le curve, verifica l'integrazione del modello nel solver; il confronto
con le curve di \hyfoam{} estratte dalle figure vettoriali del paper e con le
temperature di equilibrio ricavate dalla conservazione dell'energia ne verifica
la fisica.
